\documentclass[12pt,letterpaper]{report}

\usepackage[utf8]{inputenc}
\usepackage[T1]{fontenc}
\usepackage[spanish,es-tabla]{babel}
\usepackage{geometry}
\usepackage{graphicx}
\usepackage{xcolor}
\usepackage{hyperref}
\usepackage{longtable}
\usepackage{array}
\usepackage{booktabs}
\usepackage{tabularx}
\usepackage{enumitem}
\usepackage{float}
\usepackage{caption}
\usepackage{fancyhdr}
\usepackage{titlesec}
\usepackage{listings}
\usepackage{url}
\usepackage{multicol}
\usepackage{makecell}
\usepackage{tikz}
\usepackage{setspace}

\geometry{top=2.4cm,bottom=2.4cm,left=2.4cm,right=2.4cm}
\setlength{\parskip}{0.45em}
\setlength{\parindent}{0pt}
\renewcommand{\arraystretch}{1.22}

% Paleta de colores de HostSigchos
\definecolor{sigchosGreen}{HTML}{41C717}
\definecolor{sigchosDarkGreen}{HTML}{1CA81F}
\definecolor{sigchosDark}{HTML}{172033}
\definecolor{sigchosMuted}{HTML}{64748B}
\definecolor{sigchosLight}{HTML}{E8EEF5}
\definecolor{codeBack}{HTML}{F6F8FA}
\definecolor{codeFrame}{HTML}{D0D7DE}

\hypersetup{
  colorlinks=true,
  linkcolor=sigchosDarkGreen,
  citecolor=sigchosDarkGreen,
  urlcolor=sigchosGreen,
  pdftitle={Informe 3.2 del proyecto HostSigchos},
  pdfauthor={Andrade Danna, Lainez Ricardo, Santi Jeancarlo, Tufiño Erick}
}

\pagestyle{fancy}
\fancyhf{}
\lhead{\textcolor{sigchosMuted}{HostSigchos}}
\rhead{\textcolor{sigchosMuted}{Informe 3.2 - Avance Tercer Parcial}}
\cfoot{\thepage}
\renewcommand{\headrulewidth}{0.4pt}

\titleformat{\chapter}{\normalfont\Huge\bfseries\color{sigchosDarkGreen}}{\thechapter}{1em}{}
\titleformat{\section}{\normalfont\Large\bfseries\color{sigchosDarkGreen}}{\thesection}{1em}{}
\titleformat{\subsection}{\normalfont\large\bfseries\color{sigchosGreen}}{\thesubsection}{1em}{}

\lstdefinestyle{codeblock}{
  backgroundcolor=\color{codeBack},
  frame=single,
  rulecolor=\color{codeFrame},
  basicstyle=\ttfamily\footnotesize,
  keywordstyle=\bfseries\color{sigchosDarkGreen},
  breaklines=true,
  breakatwhitespace=true,
  columns=fullflexible,
  keepspaces=true,
  showstringspaces=false,
  tabsize=2,
  captionpos=b,
  literate={á}{{\'a}}1 {é}{{\'e}}1 {í}{{\'i}}1 {ó}{{\'o}}1 {ú}{{\'u}}1 {Á}{{\'A}}1 {É}{{\'E}}1 {Í}{{\'I}}1 {Ó}{{\'O}}1 {Ú}{{\'U}}1 {ñ}{{\~n}}1 {Ñ}{{\~N}}1
}

\newcommand{\tech}[1]{\textbf{\textcolor{sigchosDarkGreen}{#1}}}
\newcommand{\important}[1]{\textbf{\textcolor{sigchosGreen}{#1}}}
\newcommand{\pathtext}[1]{\texttt{\detokenize{#1}}}

\begin{document}

% ═══════════════ PORTADA ═══════════════
\begin{titlepage}
\thispagestyle{empty}
\newgeometry{left=2cm, right=1cm, top=2cm, bottom=2cm}
\begin{tikzpicture}[overlay, fill opacity=0.7]
    \fill[sigchosGreen] (15,-27) rectangle (18,4);
    \rotatebox{-90}{\fill[sigchosDarkGreen] (0,-9) rectangle (2,20);}
\end{tikzpicture}
\begin{spacing}{1.5}
{\huge \bfseries UNIVERSIDAD DE LAS FUERZAS ARMADAS}\\[10pt]
\centering
\end{spacing}
\vspace{1cm}
\begin{minipage}{2cm}
\hspace{0.5cm}\rotatebox{90}{\huge INGENIERIA EN SOFTWARE}
\end{minipage}\hfill
\begin{minipage}{9cm}
    % [Espacio para: Logo de la universidad o carrera]
    \centering
    \vspace{2cm}
    \textit{[Logo ESPE]}
\end{minipage}\hfill
\begin{minipage}{3cm}
\hspace{1cm}
\vspace{1.cm}
\begin{spacing}{4}
 {\color{white}\fontsize{60}{70}\selectfont\bfseries 2\\0\\2\\6}
\end{spacing}
\end{minipage}

\vspace{3cm}
\begin{flushleft}
{\LARGE
\textbf{Estudiantes:}\\
\hspace{1cm} Andrade Danna, Lainez Ricardo, \\
\hspace{1cm} Santi Jeancarlo, Tufiño Erick\\[10pt]
\textbf{Asignatura:} Desarrollo de Aplicaciones Móviles\\
\textbf{Tema:} Informe 3.2 -- Avance Tercer Parcial HostSigchos\\
\textbf{NRC:} 30738\\
\textbf{Docente:} Doris Karina Chicaiza Angamarca\\
\textbf{Fecha:} \today\\}
\end{flushleft}
\end{titlepage}
\restoregeometry

\pagenumbering{roman}
\tableofcontents
\clearpage
\pagenumbering{arabic}

\chapter*{Resumen ejecutivo}
\addcontentsline{toc}{chapter}{Resumen ejecutivo}
Este documento técnico corresponde al Informe 3.2 (Deber del tercer parcial) y describe el progreso y avance tecnológico de HostSigchos, un sistema multiplataforma de reservas hoteleras. Se detallan exhaustivamente la arquitectura aplicada (Clean Architecture), la estructura del código, las pantallas finalizadas, la implementación con Firebase y las áreas que se encuentran en consolidación, como las fases formales de pruebas unitarias (Testing).

\chapter{Arquitectura aplicada y estructura}

HostSigchos ha logrado un avance del 95\% en su fase de desarrollo funcional. La aplicación móvil se ha orquestado íntegramente bajo el estándar de \tech{Clean Architecture} con el patrón \tech{MVVM (Model-View-ViewModel)} y utilizando \tech{Provider} para la inyección y gestión del estado, garantizando la escalabilidad. Cabe destacar que el diseño, refactorización y depuración de la arquitectura fue asistido de manera eficiente por \textbf{Antigravity}, una herramienta avanzada de inteligencia artificial agentiva que contribuyó a la estandarización del código.

\begin{itemize}[label=\color{sigchosDarkGreen}$\blacksquare$]
    \item \textbf{Capa Domain:} Consta de 7 Entidades (Usuario, Hosteria, Habitacion, Reserva, Resena, Notificacion, ChatMessage) y más de 25 \pathtext{UseCases} completamente desvinculados de librerías externas.
    \item \textbf{Capa Data:} Implementación de 7 DataSources conectados a Firebase y 3 DataSources remotos para consumo de APIs REST (\tech{OpenWeatherMap}, Geocoding y \tech{Gemini AI}).
    \item \textbf{Capa Presentation:} Un conjunto robusto de 11 \pathtext{ViewModels} operando reactivamente sobre un catálogo de 19 pantallas completamente diseñadas.
\end{itemize}

\section{Estructura del proyecto (Clean Architecture)}
La raíz \pathtext{lib/} se ha segmentado estratégicamente de la siguiente forma:

\begin{tcolorbox}[colback=codeBack,colframe=codeFrame,boxrule=0.8pt,arc=6pt]
\begin{verbatim}
HostSigchos/lib/
|-- core/                     # Validadores, constantes, l10n
|-- data/                     # Acceso a Firestore, APIs, Models JSON
|-- domain/                   # Entidades Vainilla y Casos de uso core
|-- presentation/             # Pantallas UI (19 pantallas en total)
|   |-- routes/               # Tabla de enrutamiento estricto
|   |-- viewmodels/           # Gestores de estado MVVM (Provider)
|   |-- views/                # Pantallas UI reales
|   +-- widgets/              # Componentes UI reutilizables
|-- injection_container.dart  # Locator (GetIt) para inyección de dependencias
+-- main.dart                 # Inicializador App y MultiProvider
\end{verbatim}
\end{tcolorbox}

\chapter{Desarrollo y funcionalidad}

\section{Pantallas desarrolladas y Hardware}
Se completó el desarrollo de \textbf{19 pantallas funcionales} (superando ampliamente el mínimo de 4). El flujo completo del turista incluye el ingreso, catálogo, mapa con \tech{Geolocator/Compass}, asistente virtual y creación de reservas.

\begin{itemize}
    \item \textbf{Autenticación:} Login, Register, Forgot Password, Verificación (integrando \tech{local\_auth} para biometría y Firebase Google Sign-In).
    \item \textbf{Exploración:} Home (Dashboard), HosteríasList, HosteríaDetail, Mapa (integrado con \tech{flutter\_map} nativo sin costos de Google Maps).
    \item \textbf{IA y Hardware:} Chatbot interactivo (texto a voz con \tech{flutter\_tts} y grabación de audio con el micrófono del hardware).
    \item \textbf{Transaccional:} Habitaciones, CrearReserva, Checkout, Confirmación (con persistencia atómica en Firestore).
\end{itemize}

\begin{figure}[H]
  \centering
  \includegraphics[width=0.5\textwidth]{img/dashboard.png}
  \caption{Catálogo de exploración y búsqueda reactiva de hosterías.}
\end{figure}

\section{Base de Datos utilizada: Cloud Firestore}
Toda la persistencia de datos (usuarios, inventario, precios y transacciones) se maneja a través de \tech{Cloud Firestore}. Al ser NoSQL documental, se emplean id's de referencia para simular las claves foráneas y unificar las consultas en los repositorios locales. Se configuraron las \textit{Security Rules} para bloquear intrusiones.

Se implementó además \tech{Firebase Storage} para la carga y descarga masiva de los avatares (imágenes de perfil capturadas con \tech{image\_picker}) e imágenes pesadas del carrusel de las hosterías.

\section{Código principal explicado: Prevención del Overbooking}
El principal logro técnico fue resolver la asincronía en las reservas para evitar la sobreventa. Se utilizó una validación booleana estricta en el archivo \pathtext{verificar_disponibilidad_usecase.dart}:

\begin{lstlisting}[style=codeblock,caption={Lógica de validación matemática en el Caso de Uso de Dominio.}]
Future<bool> call(Habitacion habitacion, DateTime checkIn, DateTime checkOut, int cantidadSolicitada) async {
    // 1. Obtiene las reservas desde el repositorio (Data Layer)
    final reservas = await _repository.getReservasPorHabitacion(habitacion.id);

    int habitacionesOcupadas = 0;
    for (final r in reservas) {
      if (r.estaActiva) {
        // Logica matematica para cruce de fechas
        if (checkIn.isBefore(r.fechaCheckOut) && checkOut.isAfter(r.fechaCheckIn)) {
          habitacionesOcupadas += r.numHabitaciones;
        }
      }
    }
    return (habitacion.cantidadTotal - habitacionesOcupadas) >= cantidadSolicitada;
}
\end{lstlisting}

\chapter{Problemas, evidencia y pendientes}

\section{Evidencia de funcionamiento}
La aplicación móvil se conecta correctamente con el backend de Firebase en la nube y realiza el despliegue del mapa con el GPS del hardware. Las notificaciones Push (vía \tech{flutter\_local\_notifications}) operan cuando el estado de una reserva cambia en el panel web.

\begin{figure}[H]
  \centering
  \includegraphics[width=0.5\textwidth]{img/checkout_success.png}
  \caption{Flujo transaccional finalizado exitosamente.}
\end{figure}

\section{Problemas encontrados en la iteración actual}
\begin{itemize}[label=\color{red}$\blacksquare$]
    \item \textbf{Latencia de consultas complejas:} Al inicio, la pantalla de "Hosterías" colapsaba si la latencia era alta. Se solucionó añadiendo el paquete \tech{shimmer} para brindar un esqueleto de carga y no bloquear la UI principal.
    \item \textbf{Rate limit de Gemini API:} La capa del chatbot fallaba ocasionalmente con código \texttt{429}. Se debió abstraer el manejo del error devolviendo un estado \pathtext{Failure} ordenado al \pathtext{ViewModel} para notificar al usuario que reintente luego de unos segundos.
\end{itemize}

\section{Mejoras pendientes y siguientes pasos}
Si bien se han superado todos los hitos del tercer parcial, siempre existen oportunidades de mejora para el sistema en producción.

\textbf{Mejoras pendientes:}
\begin{itemize}[label=\color{orange}$\blacksquare$]
    \item \textbf{Pruebas de software formales (Testing):} La implementación exhaustiva de las pruebas automatizadas (Tests Unitarios utilizando \tech{flutter\_test} y \tech{mockito}, así como tests de integración con \tech{Appium}) se encuentra programada para la última fase del desarrollo antes de la generación del entregable final.
    \item \textbf{Firma del APK y Despliegue en Tiendas:} El paso inminente es generar las llaves criptográficas seguras (\textit{Keystore}) para firmar el APK en modo Release y enviarlo a Firebase App Distribution o a la Google Play Store.
    \item \textbf{Gestión de caché offline:} Añadir persistencia de caché local para que los turistas puedan consultar sus reservas aunque no tengan conexión a internet en ciertas zonas geográficas de montaña.
\end{itemize}

Todos los conflictos de dependencias actuales y los \textit{warnings} del linter fueron superados exitosamente, lo que deja el proyecto con una base sólida y limpia para poder afrontar la fase de pruebas formales sin inconvenientes técnicos.

\bibliographystyle{ieeetr}
\bibliography{referencias}

\end{document}
